\chapter{Distributed-parameter systems}\label{ch:distributed}

Transfer functions of systems governed by partial differential equations
are ratios of entire functions that often look multivalued. This chapter
shows how to write them as analytic pairs and computes their poles for the
heat, string, duct and beam examples of \citet{curtain2009}. The models are
implemented in \code{examples/models/distributed\_model.m}; the study is
\code{examples/distributed\_systems/run\_nonrational\_examples.m}.

\section{Removing apparent branch points}

For the heat equation the natural expressions contain $\sqrt s$. Define
\[
  C(q)=\cosh\sqrt q=\sum_{k\ge0}\frac{q^k}{(2k)!},\qquad
  H(q)=\frac{\sinh\sqrt q}{\sqrt q}=\sum_{k\ge0}\frac{q^k}{(2k+1)!},\qquad
  H(0)=1 .
\]
Both are entire: only integer powers of $q$ appear. The implementation
evaluates $H$ by its series for $|q|<10^{-4}$ and by the closed form
elsewhere, which avoids the $0/0$ at $q=0$ and does not introduce an
artificial pole. The same idea applies to the models below: each factor is
evaluated through a regularized formula and is analytic in the search
regions used here, so \code{AssumeAnalytic} is justified for both factors.
The heat, string and duct factors are entire; the Kelvin--Voigt beam factors
are analytic except at $s=-1/c_d$ (Section~\ref{sec:kelvinvoigt}), which the
regions must exclude.

\section{Heat equation: the sensor position}\label{sec:heat}

With the normalizations $L=\alpha=K_0=1$ of the paper, boundary control and
a point sensor at $x_0$, the Neumann, Dirichlet and mixed problems give
\[
  G_N=\frac{C(x_0^2s)}{s\,H(s)},\qquad
  G_D=\frac{x_0H(x_0^2s)}{H(s)},\qquad
  G_M=\frac{x_0H(x_0^2s)}{C(s)} .
\]
With $x_0=1/2$ and $-100<\Real s<2$, the Neumann model keeps only $0$ and
$-4\pi^2$: the candidates $-\pi^2$ and $-9\pi^2$ cancel. The Dirichlet
model keeps $-\pi^2$ and $-9\pi^2$ and cancels $-4\pi^2$. The mixed model
keeps $-\pi^2/4$, $-9\pi^2/4$ and $-25\pi^2/4$. With the Dirichlet model and
$x_0=1$, $G\equiv1$ and every candidate cancels. These are modes of the PDE
that the sensor cannot see (Figure~\ref{fig:heat}).

\begin{figure}[H]
\centering
\includegraphics[width=\textwidth]{heat_and_characteristic.png}
\caption{Poles retained for four sensor positions and three boundary
conditions (the Dirichlet row $x_0=1$ is empty because $G=1$). The last
panel shows the characteristic example of Chapter~\ref{ch:quick}.}
\label{fig:heat}
\end{figure}

\section{String and duct}

For the string of Section~2 of the paper, with actuator distribution
$b(x)=1$ on $[0,1/2]$ and $L=c=1$, the modal coefficients are proportional to
$1-\cos(k\pi/2)$, which vanishes for $k$ multiple of 4; those modes are not
poles. The damped model is the output feedback $G/(1+\varepsilon G)$ of the
paper, with $\varepsilon=0.5$, and the removable factor at $s=0$ is
regularized by a series.

For the acoustic duct with constant reflectance $a=0.5$ and $L=c=1$, the
poles lie on the vertical line $\ln(a)/2+\ii k\pi$. With the radiation
impedance of Section~3.4 (Fig.~5 parameters of the paper: $L=3.54$\,m,
$c=341$\,m/s, $\rho=1.20$\,kg/m$^3$, radius $0.101$\,m, sensor at $L/2$), the
reflectance depends on frequency and the poles bend to the left. Its
rational denominator is cleared from $N$ and $D$ together, so that no
artificial poles are introduced (Figure~\ref{fig:waveduct}).

Visible poles of the damped string in the left half-plane do not prove
internal stability of modes hidden from this input/output pair.

\begin{figure}[H]
\centering
\includegraphics[width=\textwidth]{wave_duct_beam.png}
\caption{String (undamped, damped, and cancelled modes), duct with
constant reflectance and with radiation impedance, and beam poles compared
with an independent modal formula.}
\label{fig:waveduct}
\end{figure}

\section{Kelvin--Voigt beam: an accumulation point}\label{sec:kelvinvoigt}

For the shear-force to tip-velocity transfer function of Section~4.2 of the
paper, with $L=EI=I=1$ and Kelvin--Voigt damping $c_d=0.02$, let
\[
  q(s)=\frac{-s^2}{1+c_ds},\qquad m=q^{1/4},\qquad
  G(s)=\frac{s\,[\cosh m\sin m-\sinh m\cos m]/m^3}{(1+c_ds)[1+\cosh m\cos m]} .
\]
The factors containing $m$ are invariant under $m\mapsto\ii m$, so the choice
of fourth root disappears except at $s=-1/c_d$; the numerator ratio has the
removable value $2/3$ at $m=0$. With the positive roots $\beta_k$ of
$1+\cosh\beta\cos\beta=0$, the poles solve $s^2+c_d\beta_k^4s+\beta_k^4=0$, an
independent reference. In $[-25,2]\times[-160,160]$ the search returns
$-0.12362363\pm3.51384128\ii$ and $-4.85518819\pm21.49292828\ii$, within
$10^{-14}$ of the reference.

As $\beta_k\to\infty$, one branch of these poles converges to $-1/c_d=-50$
(Figure~\ref{fig:accumulation}). Every neighborhood of $-50$ contains
infinitely many poles: $G$ is not meromorphic there, and no finite contour
computation can resolve a region containing that point. The model reports
it in \code{meta.singularities}, and the option \code{Singularities}
rejects any region that contains it.

\begin{figure}[H]
\centering
\includegraphics[width=\textwidth]{beam_accumulation.png}
\caption{Modal reference with 80 modes: one branch of poles converges to
$s=-1/c_d=-50$, the other moves to the left.}
\label{fig:accumulation}
\end{figure}

\section{Results}

\begin{table}[H]\centering\small
\caption{Poles in the stated rectangles; ``cancelled'' counts locations.
The reference errors compare with closed-form or modal formulas that are
never used as seeds.}\label{tab:distributed}
\begin{tabular}{lrrrrr}\toprule
Model & $[\Real_{\min},\Real_{\max}]$ & $|\Imag|<$ & Poles & Cancelled & Max. error\\\midrule
Heat, Neumann & $[-100,2]$ & 3 & 2 & 2 & $5\times10^{-14}$\\
Heat, Dirichlet & $[-100,2]$ & 3 & 2 & 1 & $1\times10^{-14}$\\
Heat, mixed & $[-100,2]$ & 3 & 3 & 0 & $7\times10^{-15}$\\
String, undamped & $[-2,1]$ & 30 & 14 & 4 & $3\times10^{-21}$\\
String, $\varepsilon=0.5$ & $[-2,1]$ & 30 & 14 & 4 & ---\\
Duct, constant reflectance & $[-2,1]$ & 20 & 13 & 0 & $2\times10^{-12}$\\
Duct, radiation impedance & $[-400,30]$ & 2500 & 16 & 0 & ---\\
Beam, Kelvin--Voigt & $[-25,2]$ & 160 & 4 & 0 & $7\times10^{-15}$\\
\bottomrule\end{tabular}
\end{table}

For the damped string and the radiating duct there is no independent
closed form; they are verified by the contour counts and the residuals of
the original equations.
